\section{Related work}
\paragraph{Lifting visual evidence to scene representations.}
VALA~\cite{vala} uses visibility-aware feature lifting and robust multi-view
aggregation for open-vocabulary Gaussian segmentation. LUDVIG~\cite{ludvig}
aggregates lifted visual features and applies graph diffusion; Occam's
LGS~\cite{occam} provides an efficient Gaussian language representation.
These approaches establish the value of correspondences, visibility and
aggregation. We keep the available teacher representations and focus on
forming query responses from them. Neither weighted averaging, robust fitting
nor frozen native features alone is claimed as our contribution.

\paragraph{Query-dependent evidence and complete extent.}
EviSplat~\cite{evisplat} preserves observation features in class-agnostic
instances and combines query-selected evidence with local Gaussian support.
Query-dependent use of preserved evidence is therefore prior art. Our focus is
the distinction between identity applicability and authority to support or
exclude spatial extent, including completion of reference scope before
transport and compatibility of response objectives with actual selection.
The distinction is evaluated through matched interventions; the formulation
does not establish that our scope estimates are universally more accurate.

\paragraph{Teacher capabilities and response-level adaptation.}
SAM3~\cite{sam3}, SigLIP2~\cite{siglip2}, DINOv2~\cite{dinov2} and
Mask-Adapter~\cite{maskadapter} supply attributed segmentation or matching
capabilities. We study how their evidence is used, not a new foundation model.
For reference queries, frozen image decoding and tracking are combined with
reference-side admission and geometric response fitting. For map-only queries,
complete membership includes independent labelled-room supervision. Published
method comparisons and retained local reproductions are contextual when
teacher access, prompts, geometry or evaluation domains differ; our component
controls provide the narrower causal evidence.
